\section{Multimodalidad: de la comprensión de video a la fusión de modelos generativos}

Esta sección aborda la cuarta línea principal: la multimodalidad. El desarrollo multimodal no consiste simplemente en «agregarle imágenes a un modelo de lenguaje», sino en la convergencia gradual de la comprensión visual, el modelado de video, los modelos generativos, el Transformer visual, la alineación imagen-texto y el Transformer de difusión. La palabra central aquí es fusión: las distintas modalidades terminan por influirse mutuamente dentro de una representación unificada, una generación unificada o un marco de razonamiento unificado.

\begin{figure}[H]
\centering
\includegraphics[width=0.96\textwidth]{multimodal-lineage.png}
\caption{Genealogía de los modelos multimodales: convergencia de video, GAN, Diffusion, ViT, CLIP, Stable Diffusion y DiT. Diagrama conceptual propio, elaborado a partir de la conversación entre 03:08:08 y 03:56:38.}
\end{figure}

\begin{knowledgebox}{Lectura de la figura: la multimodalidad no es una línea recta}
La comprensión de video, la generación de imágenes, el Transformer visual, la alineación imagen-texto y los modelos de difusión partieron de problemas distintos. Convergieron después porque los modelos necesitan comprender el mundo y generarlo a la vez, y conectar la capacidad visual con la intención del usuario por medio del lenguaje.
\end{knowledgebox}

\subsection{DeepVideo, redes de dos flujos y la comprensión temprana de video}

DeepVideo y las redes de dos flujos representan los primeros intentos del aprendizaje profundo por entrar al ámbito del video. Una vez que las imágenes pudieron procesarse con CNN, los investigadores se preguntaron con naturalidad si el video también podía tratarse así. La dificultad es que el video añade la dimensión temporal: el modelo no solo debe reconocer el contenido de cada cuadro, sino también comprender acciones, movimiento y eventos.

\begin{importantbox}{Los problemas que el video añade frente a la imagen}
El modelado de video añade al menos tres problemas: la dependencia temporal, la representación del movimiento y los eventos de largo alcance. Una imagen aislada responde «qué es esto»; el video además debe responder «cómo cambia», «por qué cambia así» y «qué ocurrirá después».
\end{importantbox}

\subsection{GAN, Diffusion y DDPM}

La línea de generación de imágenes estuvo dominada primero por las GAN. Las GAN generan muestras nítidas, pero su entrenamiento es inestable; los VAE se entrenan de forma estable, pero sus imágenes tienden a salir borrosas. Diffusion creció en sus inicios a la sombra de las GAN, y DDPM devolvió a los modelos de difusión al centro del escenario de la generación de imágenes. La ventaja de los modelos de difusión es que su entrenamiento es estable, la calidad de generación es buena y se combinan bien con el control condicional y la guía por texto.

\begin{knowledgebox}{Ventajas y desventajas de GAN frente a Diffusion}
\begin{tabularx}{\textwidth}{p{0.22\textwidth}p{0.34\textwidth}X}
\toprule
Familia de modelos & Ventajas & Dificultades \\
\midrule
GAN & Muestras nítidas, generación rápida & Entrenamiento inestable, alto riesgo de colapso de modos. \\
VAE & Entrenamiento estable, interpretación como modelado probabilístico & Las muestras tienden a salir borrosas. \\
Diffusion & Estable, de alta calidad, fácil de controlar condicionalmente & Costo de muestreo elevado; requiere optimización de ingeniería. \\
\bottomrule
\end{tabularx}
\end{knowledgebox}

\subsection{ViT, CLIP, Stable Diffusion y DiT}

La subsección anterior mostró cómo los modelos generativos avanzaron hacia la difusión; esta examina cómo la visión y el lenguaje convergen en una arquitectura unificada. ViT divide la imagen en patches y procesa los tokens visuales con un Transformer; CLIP coloca imágenes y textos en un mismo espacio semántico y se convirtió en un pilar de la generación de imágenes a partir de texto y de la búsqueda multimodal; Stable Diffusion combinó los modelos de difusión, el espacio latente y el ecosistema de código abierto, e impulsó la popularización de la generación de imágenes; DiT, por su parte, fusiona más a fondo Diffusion y Transformer, y representa la expectativa de una futura arquitectura unificada.

\begin{importantbox}{Por qué CLIP es fundamental}
La importancia de CLIP reside en la alineación imagen-texto. Permite que el modelo sepa si una imagen y un texto coinciden semánticamente, lo que sustenta la generación de imágenes a partir de texto, la búsqueda de imágenes, el filtrado de datos y el aprendizaje de representaciones multimodales. Sin alineación imagen-texto, a un modelo generativo le resulta muy difícil entender de manera estable las indicaciones del usuario.
\end{importantbox}

\subsection{Resumen de la sección}

La línea multimodal muestra que la visión, el lenguaje y la generación se están acercando. Las GAN establecieron el punto de partida de la generación, Diffusion resolvió la calidad y la estabilidad, ViT llevó la visión al Transformer, CLIP tendió el puente entre imagen y texto, y DiT terminó de fusionar la generación con el Transformer.

